\section{Worked Example: Stepping Through the Calculus}\label{sec:worked_example}\label{sec:filter_example}

This appendix traces the full execution of the Filtered Stepper
Calculus on a small example, showing each phase of the stepping
pipeline in detail. Readers may find it helpful to follow this
example while reading the formal rules in Section~\ref{sec:filter}.

Consider the following Hazel program:
\lstset{xleftmargin=1in}
\begin{lstlisting}[language=hazel,caption={A simple Hazel program}]
debug eval(1 + 2 + 3 + 4) in
debug stop(3 + 3 + 4) in
1 + 2 + 3 + 4\end{lstlisting}

This program sums the numbers from 1 to 4. The plus operation in Hazel
is binary and left-associative, so the expression
\code{1 + 2 + 3 + 4} is evaluated as follows:
\code{1 + 2} becomes
\code{3}, then \code{3 + 3}
becomes \code{6}, and finally
\code{6 + 4} becomes
\code{10}.

Each ``macroscopic'' step of the evaluation is achieved by a series of
``microscopic'' steps. Here is a detailed description of the first few
``microscopic'' steps of the program above.

\begin{enumerate}
\item \emph{Instrumentation}. \label{num:simple_example_instrument} In this
  step, we traverse the expression recursively, collecting all filters along the
  path. For each sub-expression, we check if it matches any of the collected
  filters. If a match is found, we insert a residue expression with the filter
  information around the matched sub-expression.

  The matching process itself is recursive. If an
  expression is constructed in the same way as the filter pattern, we recurse into
  their sub-expressions to check if the sub-expressions match. If all
  sub-expressions match, we say the expression matches the filter. In
  most cases, we want the filter pattern to match
  against more kinds of expression. Therefore, we introduced two
  wildcard matchers: \code{\$e} and
  \code{\$v}. The former matches against any
  expression, while the latter matches against any value. Assume
  we are checking whether expression \(e\) matches filter
  \(f\). This process can be informally described as the following
  steps.

  \begin{enumerate}
  \item If \(f\) is the wildcard \(\$e\), exit the process and return
    that \(e\) and \(\$e\) match.
  \item If \(f\) is the wildcard \(\$v\), and \(e\) is a value, exit the process and return
    that \(e\) matches \(\$v\).
  \item If \(f\) is the wildcard \(\$v\), and \(e\) is not a value, exit
    the process and return that \(e\) does not match \(\$v\).
  \item If \(e\) and \(f\) do not share a common construction,
    exit the matching process and return that \(e\) and \(f\) do not
    match.
  \item Otherwise, iterate over each pair of child expressions or
    child filters in \(e\) and \(f\). Since \(e\) and \(f\) share the
    same construction, their number of child items is always the
    same. For each pair, go to step 1.
  \item If every pair in step 5 matches, exit the process and return that \(e\) and \(f\) match.
  \end{enumerate}

  If no sub-expression matches any of the collected filters, the
  stepper defaults to stepping through every expression.

  In our example, we apply three filters to the expression
  \code{1 + 2 + 3 + 4}:
  \begin{enumerate}
  \item \code{debug stop(\$e)}
  \item \code{debug eval(1 + 2 + 3 + 4)}
  \item \code{debug stop(3 + 3 + 4)}
  \end{enumerate}

  Before instrumentation, we assign priorities to active
  filters. The priority is assigned based on the order in which filters are added to
  the list of active filters. Here, we assign priority 0 to the first
  filter, 1 to the second one, and 2 to the last one. The priority
  is used when deciding which action to take for the next step.

  During the instrumentation of the first filter, any sub-expression
  will match the pattern, so we wrap every sub-expression in a residue
  expression \(\Residue{\ActStep}{\GasOne}{0}{\ldots}\). For the
  innermost sub-expressions, like \code{1},
  \code{2}, etc., we do not instrument them since
  they are values. Because a value cannot be evaluated, it would make
  no sense if we allow the stepper to stop before ``evaluating'' a
  value.

  During the instrumentation of the second filter, the sub-expression
  \code{1 + 2 + 3 + 4} matches the pattern, so we
  wrap it in a residue expression
  \(\Residue{\ActSkip}{\GasAll}{0}{\ldots}\) (I-Add-Y).  The
  instrumented expression is shown in
  Listing~\ref{lst:simple_example_instrument}.

  \begin{lstlisting}[language=hazel,caption={Program after instrumentation},label={lst:simple_example_instrument}]
debug stop(\$e) in
debug eval(1 + 2 + 3 + 4) in
debug stop(3 + 3 + 4) in
@\(
\prescript{\ActStep}{0}{\langle}
\prescript{\ActSkip}{1}{\langle}
\prescript{\ActStep}{0}{\langle}
\prescript{\ActStep}{0}{\langle}
\)@1 + 2@\({\rangle}^{\GasOne}\)@ + 3@\({\rangle}^{\GasOne}\)@ + 4@\(\rangle^{\GasAll}
\rangle^{\GasOne}\)@\end{lstlisting}

  One thing you might have noticed is that we did not put instrumented
  residue expressions around values.

  During the instrumentation of the third filter, no sub-expression matches the
  pattern \code{3 + 3 + 4}, so no residue expression is
  inserted (I-Add-N). The instrumented expression remains unchanged.

\item \emph{Optimization}. The stepper might have inserted too many
  residue expressions in the instrumentation step, and this could
  cause a significant slowdown of the stepper. Therefore,
  some redundant residue expressions are optimized out during this
  step.

  Generally, we want to remove all immediately adjacent residue
  expressions by picking the one with highest priority, because this is
  the only one that can be picked among these residue expressions
  during the action selection process. Here, the two outermost
  residue expressions are adjacent:
  \(\Residue{\ActStep}{\GasOne}{0}{\Residue{\ActSkip}{\GasAll}{1}{\ldots}}\),
  we remove the ones with lower priorities (O-Residue-Inner) to get
  \(\Residue{\ActSkip}{\GasAll}{1}{\ldots}\).

\begin{lstlisting}[
  language=hazel,
  caption={Program after instrumentation, optimized},
]
debug stop(\$e) in
debug eval(1 + 2 + 3 + 4) in
debug stop(3 + 3 + 4) in
@\(
\prescript{\ActSkip}{1}{\langle}
\prescript{\ActStep}{0}{\langle}
\prescript{\ActStep}{0}{\langle}
\)@1 + 2@\({\rangle}^{\GasOne}\)@ + 3@\({\rangle}^{\GasOne}\)@ + 4@\(
\rangle^{\GasAll}
\)@\end{lstlisting}

\item \emph{Decomposition}. \label{num:simple_example_decompose} The
  expression is then decomposed into an evaluation context and a
  selected expression. The instrumented expression from the last step can
  be decomposed by recursively decomposing addition (D-Add-L) and
  residue expressions (D-Residue) into an evaluation context
  (Listing~\ref{lst:simple_example_decompose}) and a selected
  expression \code{1 + 2}.

  \begin{lstlisting}[language=hazel,caption={Decomposed evaluation context},label={lst:simple_example_decompose}]
debug stop(\$e) in
debug eval(1 + 2 + 3 + 4) in
debug stop(3 + 3 + 4) in
@\(
\prescript{\ActSkip}{1}{\langle}
\prescript{\ActStep}{0}{\langle}
\prescript{\ActStep}{0}{\langle}
\circ{\rangle}^{\GasOne}\)@ + 3@\({\rangle}^{\GasOne}\)@ + 4@\(
\rangle^{\GasAll}
\)@\end{lstlisting}

\item \emph{Making a Decision}. \label{num:simple_example_select} Decide
  which action (stop or continue) to take based on all the residue
  expressions around the mark in the evaluation context. The action is
  decided using the following set of informal rules:

  \begin{enumerate}
  \item For an expression that is not a residue expression, the action
    will be the one that comes from the sub-expression that contains
    the mark.
  \item For a residue expression, if the action from its only
    sub-expression has a higher priority, use the action from the
    sub-expression; otherwise, use the action of the residue
    expression.
  \end{enumerate}

  In this case, the evaluation context yields an action of
  \(\ActSkip\) in the decision-making process, while the expression
  itself is kept the same.

\item \emph{Residue removal (Decaying)}. Residue expressions that
  have a gas of one are removed from the expression. This step is
  done using the \(\Decay\) function. The expression after the removal
  of the residue expression is:
  \begin{lstlisting}[language=hazel,caption={Evaluation context after decaying},label={lst:simple_example_decay}]
debug stop($e) in
debug eval(1 + 2 + 3 + 4) in
debug stop(3 + 3 + 4) in
@\(\prescript{\ActSkip}{1}{\langle}\circ\)@ + 3 + 4@\(\rangle^{\GasAll}\)@\end{lstlisting}

\item \emph{User Interaction}. If the selected expression produced in
  Step~\ref{num:simple_example_decompose} is of the form
  \code{debug .. in ..} or
  \(\Residue{a}{g}{l}{e}\), or the selected action from
  Step~\ref{num:simple_example_select} is \(\ActSkip\), the stepper
  will skip over this step and proceed into the next step. (S-Filter and
  S-Skip). Otherwise, the stepper will yield the control to the user,
  and let them choose which next possible step to take. (S-Step)

  In this case, the selected action is \(\ActSkip\). Therefore, the stepper
  will skip over this step and proceed into the next step.

\item \emph{Instruction Transition}. The selected expression is
  instruction-transitioned to a new expression. In this case, the selected
  expression \code{1 + 2} is
  instruction-transitioned to \code{3} according
  to the transition rule (T-Add).

\item \emph{Composition}. When the selected expression is
  instruction-transitioned to a new expression, we need to compose the
  transitioned selected expression with the evaluation context to get the
  updated expression. The rules for composition are the same as the
  rules for decomposition, but instead, we start with rule (D-Add-L)
  to build up the recomposed expression in a bottom-up way. The
  recomposed expression is:
  \begin{lstlisting}[language=hazel]
debug stop($e) in
debug eval(1 + 2 + 3 + 4) in
debug stop(3 + 3 + 4) in
@\(\prescript{\ActSkip}{1}{\langle}\)@3 + 3 + 4@\(\rangle^{\GasAll}\)@\end{lstlisting}

\item \emph{Recursion}. As long as we do not get a value, we will keep
  stepping through the expression, i.e. go back to
  Step~\ref{num:simple_example_instrument}.
\end{enumerate}

In the following round of the procedure, the program is first
instrumented (Listing~\ref{lst:simple_example_1_instrument}), then is
decomposed into an evaluation context
(Listing~\ref{lst:simple_example_1_decompose}) with a selected
expression \code{3 + 3}.

\begin{lstlisting}[
  language=hazel,
  caption={Instrumented program in Round 2},
  label={lst:simple_example_1_instrument}
]
debug stop(\$e) in
debug eval(1 + 2 + 3 + 4) in
debug stop(3 + 3 + 4) in
@\(
\prescript{\ActSkip}{1}{\langle}
\prescript{\ActStep}{0}{\langle}
\prescript{\ActStep}{2}{\langle}
\prescript{\ActStep}{0}{\langle}
\)@3 + 3@\(
\rangle^{\GasOne}
\)@ + 4@\(
\rangle^{\GasOne}
\rangle^{\GasOne}
\rangle^{\GasAll}
\)@\end{lstlisting}

\begin{lstlisting}[
  language=hazel,
  caption={Decomposed evaluation context in Round 2},
  label={lst:simple_example_1_decompose}
]
debug stop(\$e) in
debug eval(1 + 2 + 3 + 4) in
debug stop(3 + 3 + 4) in
@\(
\prescript{\ActSkip}{1}{\langle}
\prescript{\ActStep}{0}{\langle}
\prescript{\ActStep}{2}{\langle}
\prescript{\ActStep}{0}{\langle}
\circ
\rangle^{\GasOne}
\)@ + 4@\(
\rangle^{\GasOne}
\rangle^{\GasOne}
\rangle^{\GasAll}
\)@\end{lstlisting}

Then, we compute the selected action to be \(\ActStep\). Therefore,
the stepper stops at this step and lets the user choose which
sub-expression should be evaluated for the next step by clicking the
yellow area shown in Figure~\ref{fig:simple_example_final_ui}.

\begin{figure}[h]
  \centering
  \frame{\includegraphics[width=0.4\textwidth]{images/filter-example-user-select.png}}
  \Description{Stepper UI showing the expression 1 + 2 + 3 + 4 with a highlighted subexpression selection and control buttons.}
  \caption{User Interface for the \emph{User Selection} Step.}\label{fig:simple_example_final_ui}
\end{figure}

%%% Local Variables:
%%% mode: LaTeX
%%% TeX-master: "main"
%%% End:
